% =============================================================================
% exercises/ejercicios_cap03.tex
% Ejercicios propuestos — Capítulo 3: Velocidad y Equilibrio
% =============================================================================

\section*{Ejercicios Propuestos}
\addcontentsline{toc}{section}{Ejercicios propuestos}
\label{sec:ejercicios_cap03}

\begin{bloque_ejercicios}[Ejercicios --- Cap\'{\i}tulo 3: Velocidad y Equilibrio]

\begin{ejercicios}

% ----- Ejercicio 1: Teoría de colisiones y energía de activación -----------
\item La teor\'{\i}a de las colisiones establece que s\'olo los choques
con energ\'{\i}a superior a la \textbf{energ\'{\i}a de activaci\'on} $E^*$ y
con la \textbf{orientaci\'on adecuada} producen reacci\'on.

\begin{enumerate}[label=\alph*)]
    \item Para la reacci\'on elemental bimolecular
          \ce{A + B -> C + D}, la constante de velocidad
          seg\'un la teor\'{\i}a de colisiones es $k = Z'\,e^{-E/RT}$.
          Identifique qu\'e representa cada factor ($Z'$ y $e^{-E/RT}$)
          y explique c\'omo afecta cada uno a la velocidad cuando se
          incrementa la temperatura.

    \item Con base en el perfil de energ\'{\i}a de la
          \textbf{Figura~\ref{e-activa}}, indique cu\'al de las
          dos reacciones (exot\'ermica o endot\'ermica) tiene un
          mayor $|\Delta H|$ y cu\'al requiere mayor energ\'{\i}a de
          activaci\'on para la reacci\'on \textbf{inversa}.
          Justifique gr\'aficamente.

    \item La combusti\'on del magnesio es altamente exot\'ermica y
          sin embargo requiere ser calentado a alta temperatura para
          iniciar. Explique esta aparente contradicci\'on en t\'erminos
          de la energ\'{\i}a de activaci\'on y el complejo activado.

    \item El factor de frecuencia $A$ para reacciones bimoleculares
          es del orden de $10^9$--$10^{10}$ \si{\mole\per\liter} a
          \SI{300}{\kelvin}. Si $k = Z'\,e^{-E/RT}$ y
          $E = \SI{80}{\kilo\joule\per\mole}$, calcule la relaci\'on
          $k/A = e^{-E/RT}$ a \SI{300}{\kelvin}.
          Dato: $R = \SI{8.314}{\joule\per\mole\per\kelvin}$.
\end{enumerate}

\textit{[Respuesta parcial:
(d) $e^{-E/RT} = e^{-80000/(8.314 \times 300)}
    = e^{-32.09} \approx 9.1\times10^{-15}$;
    solo una fracción muy peque\~na de colisiones es reactiva.]}

% ----- Ejercicio 2: Factores que influyen en la velocidad ------------------
\item Para cada situaci\'on, identifique el \textbf{factor de velocidad}
involucrado (concentraci\'on, temperatura, superficie de contacto,
catalizador o naturaleza de los reactivos) y prediga si la velocidad
aumenta o disminuye:

{%
\centering
\begin{tabular}{clp{3cm}c}
\toprule
\textbf{N\'um.} & \textbf{Situaci\'on} & \textbf{Factor} & \textbf{Efecto} \\
\midrule
a) & Refrigerar alimentos para conservarlos & & \\
b) & Moler un mineral antes de lixiviarlo   & & \\
c) & Usar \ce{MnO2} en la descomposici\'on
      de \ce{H2O2}                           & & \\
d) & Duplicar la concentraci\'on de \ce{HCl}
      en la reacci\'on con Zn                & & \\
e) & Usar \ce{NaCl} disuelto en lugar de
      \ce{NaCl} s\'olido                     & & \\
f) & Aumentar la presi\'on en la s\'{\i}ntesis
      del amon\'{\i}aco (fase gaseosa)           & & \\
\bottomrule
\end{tabular}
\par
}%

\medskip
Escriba adem\'as la reacci\'on qu\'{\i}mica para los incisos
\textbf{(c)} y \textbf{(d)}, tomando como referencia los ejemplos
del cap\'{\i}tulo.

% ----- Ejercicio 3: Constante de equilibrio ---------------------------------
\item Escriba la expresi\'on de la constante de equilibrio $K_{eq}$
para cada una de las siguientes reacciones y clasif\'{\i}quelas seg\'un el
valor esperado de $K_{eq}$ (muy grande $\gg1$, intermedio $\approx1$
o muy peque\~no $\ll1$) con base en la informaci\'on del cap\'{\i}tulo:

\begin{enumerate}[label=\alph*)]
    \item \ce{3H2_{(g)} + N2_{(g)} <=> 2NH3_{(g)}}
    \item \ce{CO_{(g)} + 2H2_{(g)} <=> CH3OH_{(g)}}
    \item \ce{H2_{(g)} + I2_{(g)} <=> 2HI_{(g)}}
          \hfill ($K_{eq} = 54.8$ a $425^\circ$C)
    \item \ce{COCl2_{(g)} <=> CO_{(g)} + Cl2_{(g)}}
          \hfill ($K_{eq} = 7.6\times10^{-4}$ a $400^\circ$C)
    \item \ce{2NO2_{(g)} <=> N2O4_{(l)}}
    \item \ce{2H2O_{(l)} <=> 2H2_{(g)} + O2_{(g)}}
\end{enumerate}

Para los incisos \textbf{(c)} y \textbf{(d)}, indique adem\'as
si el equilibrio se desplaza hacia los productos o hacia los reactivos
y qu\'e implica esto para la estabilidad del compuesto.

% ----- Ejercicio 4: Principio de Le Chatelier ------------------------------
\item La s\'{\i}ntesis industrial del amon\'{\i}aco se realiza mediante el
\textbf{proceso Haber}:
\[
    \ce{N2_{(g)} + 3H2_{(g)} <=> 2NH3_{(g)}} \quad \Delta H = -92\,\text{kJ/mol}
\]

Para cada perturbaci\'on al sistema en equilibrio, prediga usando el
\textbf{principio de Le Chatelier} la direcci\'on del desplazamiento
(hacia productos o reactivos) y justifique:

{%
\centering
\begin{tabular}{clp{3.5cm}p{3cm}}
\toprule
 & \textbf{Perturbaci\'on} & \textbf{Desplazamiento} & \textbf{Raz\'on} \\
\midrule
a) & Se aumenta $[\ce{N2}]$       & & \\
b) & Se aumenta la temperatura    & & \\
c) & Se extrae \ce{NH3} del sistema & & \\
d) & Se aumenta la presi\'on total & & \\
e) & Se agrega un catalizador      & & \\
f) & Se disminuye $[\ce{H2}]$      & & \\
\bottomrule
\end{tabular}
\par
}%

\medskip
?`Por qu\'e industrialmente se opera a \textbf{alta presi\'on} y
\textbf{temperatura moderada} a pesar de que la reacci\'on es
exot\'ermica? Explique el compromiso cin\'etico-termodin\'amico.

% ----- Ejercicio 5: Brønsted-Lowry, pares conjugados y anfóteros -----------
\item Con base en la teor\'{\i}a de Br\"onsted-Lowry:

\begin{enumerate}[label=\alph*)]
    \item Para cada reacci\'on, identifique el \textbf{\'acido},
          la \textbf{base}, el \textbf{\'acido conjugado} y la
          \textbf{base conjugada}:
          \begin{enumerate}[label=\roman*)]
              \item \ce{HCl_{(g)} + H2O_{(l)} -> H3O+_{(ac)} + Cl^-_{(ac)}}
              \item \ce{NH3_{(ac)} + H2O_{(l)} <=> NH4+_{(ac)} + OH^-_{(ac)}}
              \item \ce{HCO3^-_{(ac)} + H2O_{(l)} <=>
                        H2CO3_{(ac)} + OH^-_{(ac)}}
          \end{enumerate}

    \item El \ce{Zn(OH)2} es un compuesto anf\'otero. Escriba las
          dos reacciones que demuestran este comportamiento
          (una con \ce{HCl} y otra con \ce{NaOH}), identificando
          en cada caso si act\'ua como \'acido o como base.

    \item Clasifique los siguientes compuestos como \'acido fuerte,
          \'acido d\'ebil, base fuerte o base d\'ebil, con base en el
          \textbf{Cuadro~\ref{tabla9}} y los ejemplos del cap\'{\i}tulo:
          \ce{HClO4}, \ce{NH4+}, \ce{NaOH}, \ce{HCOOH}, \ce{Ca(OH)2},
          \ce{H3PO4}.

    \item Escriba las reacciones de neutralizaci\'on de:
          \begin{enumerate}[label=\roman*)]
              \item \ce{2HNO3 (ac) + Ca(OH)2 (ac) ->}
              \item \ce{H2SO4 (ac) + MgO (s) ->}
              \item \ce{2HCl (ac) + Na2CO3 (ac) ->}
          \end{enumerate}
          Identifique la sal formada en cada caso.
\end{enumerate}

% ----- Ejercicio 6: Cálculo de pH y pOH ------------------------------------
\item Realice los siguientes c\'alculos de $pH$ y $pOH$, usando las
relaciones $pH = -\log[\ce{H+}]$ y $pH + pOH = 14$
(\textbf{Cuadro~\ref{ph-poh}}):

\begin{enumerate}[label=\alph*)]
    \item Calcule el $pH$ y el $pOH$ para las siguientes
          concentraciones de \ce{H+}:

{%
\centering
\begin{tabular}{lcccc}
\toprule
\textbf{Soluci\'on} & $[\ce{H+}]$ (mol/L) & $pH$ & $pOH$ &
\textbf{Caracter} \\
\midrule
\'Acido closh\'{\i}drico 0.1 M & $1\times10^{-1}$ & & & \\
Vinagre            & $2\times10^{-3}$  & & & \\
Sangre             & $3.98\times10^{-8}$ & & & \\
Agua de mar        & $6.3\times10^{-9}$ & & & \\
NaOH 0.1 N        & $1\times10^{-13}$  & & & \\
\bottomrule
\end{tabular}
\par
}%

    \item ?`Cu\'al es la $[\ce{H+}]$ en el jugo de lim\'on
          ($pH = 2.2$) y en la saliva ($pH = 6.5$)?
          ?`Cu\'antas veces m\'as \'acido es el jugo de lim\'on
          que la saliva?

    \item La bilis tiene $pH = 7.0$ y la leche de magnesia
          $pH = 10.5$. Calcule $[\ce{OH^-}]$ para cada una
          usando $K_w = 1\times10^{-14}$ y verifique que
          $pH + pOH = 14$.
\end{enumerate}

\textit{[Respuesta parcial:
(b) $[\ce{H+}]_{\text{lim\'on}} = 10^{-2.2}
    = 6.3\times10^{-3}$~mol/L;
    $[\ce{H+}]_{\text{saliva}} = 10^{-6.5}
    = 3.2\times10^{-7}$~mol/L;
    raz\'on $= 6.3\times10^{-3}/3.2\times10^{-7}
    \approx 19\,700$ veces]}

% ----- Ejercicio 7: Henderson-Hasselbalch y soluciones buffer ---------------
\item La ecuaci\'on de Henderson-Hasselbalch permite calcular el $pH$
de soluciones amortiguadoras (\textit{buffer}):
\[
    pH = pK_a + \log\frac{[\text{base conjugada}]}{[\text{\'acido d\'ebil}]}
\]

\begin{enumerate}[label=\alph*)]
    \item Verifique el ejemplo del cap\'{\i}tulo: calcule el $pH$ de una
          mezcla de \'acido ac\'etico \SI{0.2}{\mole\per\liter} y acetato de
          sodio \SI{0.3}{\mole\per\liter} con $pK_a = 4.76$.

    \item Se prepara una soluci\'on amortiguadora de \'acido f\'ormico
          (HCOOH) y formiato de sodio (\ce{HCOONa}).
          Usando $K_a = 1.8\times10^{-4}$ del
          \textbf{Cuadro~\ref{tabla9}}:
          \begin{enumerate}[label=\roman*)]
              \item Calcule $pK_a$.
              \item Calcule el $pH$ si $[\ce{HCOOH}] = \SI{0.1}{\mole\per\liter}$
                    y $[\ce{HCOONa}] = \SI{0.4}{\mole\per\liter}$.
          \end{enumerate}

    \item Para una soluci\'on de ion amonio/amon\'{\i}aco con
          $K_a(\ce{NH4+}) = 5.6\times10^{-10}$:
          \begin{enumerate}[label=\roman*)]
              \item ?`Qu\'e relaci\'on $[\ce{NH3}]/[\ce{NH4+}]$
                    se necesita para obtener $pH = 9.0$?
              \item Si se desea preparar \SI{1}{\liter} de este
                    \textit{buffer} con una concentraci\'on total
                    de nitr\'ogeno de \SI{0.5}{\mole\per\liter}, calcule
                    los moles de \ce{NH3} y de \ce{NH4Cl}
                    necesarios.
          \end{enumerate}

    \item Explique por qu\'e la sangre humana ($pH = 7.4$) es un
          sistema amortiguador cr\'{\i}tico y cu\'al ser\'{\i}a la
          consecuencia fisiológica de que su $pH$ bajara a 7.0
          o subiera a 7.8.
\end{enumerate}

\textit{[Respuesta parcial:
(a) $pH = 4.76 + \log(0.3/0.2) = 4.76 + 0.18 = 4.94$~$\checkmark$;
(b-i) $pK_a = -\log(1.8\times10^{-4}) = 3.74$;
(b-ii) $pH = 3.74 + \log(0.4/0.1) = 3.74 + 0.60 = 4.34$;
(c-i) $\log([\ce{NH3}]/[\ce{NH4+}]) = 9.0 - 9.25 = -0.25$;
      $[\ce{NH3}]/[\ce{NH4+}] = 10^{-0.25} \approx 0.56$]}

% ----- Ejercicio 8: Reacciones de ácidos y bases ---------------------------
\item Complete y balancea las siguientes reacciones, identifique el
tipo (metal--\'acido, neutralizaci\'on, \'oxido--\'acido o
carbonato--\'acido) y nombre la sal formada:

\begin{enumerate}[label=\alph*)]
    \item \ce{Mg (s) + H2SO4 (ac) ->}
    \item \ce{3Ca(OH)2 (ac) + 2FeCl3 (ac) ->}
    \item \ce{6HCl (ac) + Fe2O3 (s) ->}
    \item \ce{H2SO4 (ac) + MgCO3 (ac) ->}
    \item \ce{2KOH (ac) + CuSO4 (ac) ->}
    \item \ce{3Zn (s) + 8HNO3 (dil) ->}
\end{enumerate}

Para el inciso \textbf{(f)}, explique por qu\'e se forma \ce{NO}
en lugar de \ce{H2}, a diferencia de las reacciones t\'{\i}picas
de \'acido con metal.

% ----- Ejercicio 9: Integrador — cinética, equilibrio y pH -----------------
\item La descomposici\'on del per\'oxido de hidr\'ogeno ocurre seg\'un:
\[
    \ce{2H2O2 (ac) ->[\ce{MnO2}] 2H2O (l) + O2 (g)}
\]

\begin{enumerate}[label=\alph*)]
    \item Clasifique el papel del \ce{MnO2} en esta reacci\'on
          e indique c\'omo modifica la energ\'{\i}a de activaci\'on
          sin alterar el $\Delta H$ de la reacci\'on.

    \item La regla emp\'{\i}rica del cap\'{\i}tulo indica que por cada
          $10^\circ$C de aumento la velocidad se duplica.
          Si a \SI{25}{\celsius} una muestra de \ce{H2O2}
          tarda 60 minutos en descomponerse a la mitad,
          ?`cu\'anto tardar\'a a \SI{55}{\celsius}?

    \item El \ce{H2O2} comercial tiene $pH \approx 4.5$.
          Calcule $[\ce{H+}]$ y $[\ce{OH^-}]$ para esta soluci\'on
          y determine si es \'acida, b\'asica o neutra.

    \item Si se agrega \ce{NaOH} a la soluci\'on de \ce{H2O2},
          ?`qu\'e predice el principio de Le Chatelier sobre
          la estabilidad del \ce{H2O2} sabiendo que su
          descomposici\'on se favorece en medio b\'asico?
\end{enumerate}

\textit{[Respuesta parcial:
(b) En 3 intervalos de $10^\circ$C (25$\to$35$\to$45$\to$55$^\circ$C)
    la velocidad se multiplica por $2^3=8$; el tiempo se reduce a
    $60/8 = 7.5$~min;
(c) $[\ce{H+}] = 10^{-4.5} = 3.16\times10^{-5}$~mol/L;
    $[\ce{OH^-}] = 10^{-9.5} = 3.16\times10^{-10}$~mol/L; \'acida]}

% ----- Ejercicio 10 (avanzado): Arrhenius y Kw a temperatura ---------------
\item \textbf{[Avanzado]} La ecuaci\'on de Arrhenius relaciona la
constante de velocidad con la temperatura:
\[
    k = A\,e^{-E_a/RT}
\]
y en su forma lineal:
\[
    \ln\frac{k_2}{k_1} = -\frac{E_a}{R}\left(\frac{1}{T_2}
    - \frac{1}{T_1}\right)
\]

\begin{enumerate}[label=\alph*)]
    \item Una reacci\'on tiene $k_1 = 2.0\times10^{-3}$~s$^{-1}$
          a $T_1 = \SI{300}{\kelvin}$ y $k_2 = 8.0\times10^{-3}$~s$^{-1}$
          a $T_2 = \SI{320}{\kelvin}$. Calcule la energ\'{\i}a de
          activaci\'on $E_a$ en \si{\kilo\joule\per\mole}.

    \item Con el valor de $E_a$ obtenido, estime $k$ a
          \SI{350}{\kelvin}. Comente si este resultado es
          consistente con la regla emp\'{\i}rica de duplicaci\'on
          de velocidad por cada $10^\circ$C.

    \item $K_w$ del agua var\'{\i}a con la temperatura:
          a \SI{25}{\celsius}, $K_w = 1.0\times10^{-14}$;
          a \SI{37}{\celsius} (temperatura corporal),
          $K_w = 2.4\times10^{-14}$.
          \begin{enumerate}[label=\roman*)]
              \item Calcule el $pH$ del agua pura
                    a \SI{37}{\celsius}.
              \item ?`Es este valor igual, mayor o menor que 7?
                    ?`Significa que el agua a \SI{37}{\celsius}
                    es \'acida? Justifique.
          \end{enumerate}

    \item Usando la ecuaci\'on de van't Hoff
          $\ln(K_2/K_1) = -(\Delta H/R)(1/T_2 - 1/T_1)$,
          estime $\Delta H$ para la autoionizaci\'on del agua
          a partir de los valores de $K_w$ a \SI{25}{\celsius}
          y \SI{37}{\celsius}. ?`Es la ionizaci\'on del agua
          endot\'ermica o exot\'ermica?
\end{enumerate}

\textit{[Respuesta parcial:
(a) $\ln(k_2/k_1) = \ln 4 = 1.386$;
    $E_a = 1.386 \times 8.314 /
    (1/300 - 1/320) = 1.386\times8.314/
    (2.083\times10^{-4}) \approx \SI{55.4}{\kilo\joule\per\mole}$;
(c-i) $[\ce{H+}] = \sqrt{2.4\times10^{-14}} = 1.55\times10^{-7}$;
      $pH = 6.81$; menor que 7 pero el agua sigue siendo neutra
      porque $[\ce{H+}]=[\ce{OH-}]$;
(d) $\Delta H > 0$, ionizaci\'on endot\'ermica
    ($\Delta H \approx +55$~kJ/mol)]}

\end{ejercicios}

\end{bloque_ejercicios}
% =============================================================================
% FIN DE ejercicios_cap03.tex
% =============================================================================
